\chapter{Patterns: form ergonomics}
\label{ch:formergo}

\begin{aside}
Forms are where the user gives data and the app gives
feedback. Good ergonomics: clear labels, sensible validation,
quick correction.
\end{aside}

\section{Labels and placeholders}

A label outside the field; a placeholder inside as a hint.
Placeholders disappear on typing; labels stay.

\section{Validation timing}

\begin{itemize}[leftmargin=*]
  \item On blur: validate when leaving the field.
  \item On submit: validate everything at once.
  \item On change: real-time, but only after first blur.
\end{itemize}

The on-blur-then-on-change pattern is the most user-friendly:
no premature errors; live feedback after first attempt.

\section{Error display}

Below the field, in red. Concise, actionable. ``Email
required'' beats ``This field is required''.

\section{Submit affordance}

The submit button is the goal. Disable it (greyed,
non-clickable) until valid. Or always enable but show
errors on attempted submit.

\section{Field ordering}

Logical groups: name first, then contact, then preferences.
Tab order matches reading order.

\section{Required vs optional}

Mark required fields, not optional. Less visual noise.

\section{Recap}

\begin{recap}
\begin{itemize}[leftmargin=*]
  \item Labels outside; placeholders inside.
  \item Validate on blur and on change.
  \item Errors are concise.
  \item Submit button reflects state.
  \item Mark required, not optional.
\end{itemize}
\end{recap}

\section*{Workshop}

\begin{enumerate}[leftmargin=*]
  \item Audit your forms; apply the timing pattern.
  \item Shorten any verbose error messages.
  \item Mark required fields only.
\end{enumerate}
